% ============================================================
%  Setmana 32 — Unitat 6 · Angles i distàncies a l'espai
%  Exercicis: 74, 76, 86, 88
% ============================================================

\setmana{32}{Unitat 6 \textperiodcentered{} Angles i distàncies a l'espai}{25 d'abril de 2027}%
        {74, 76, 86 i 88}

\begin{errada}{Exercici 74}
  \ondiu{«c) Trobem la projecció ortogonal del punt sobre la recta.»}
  \hadedir{«b) Trobem la projecció ortogonal del punt sobre la recta.»}
  \motiu{L'exercici només té dos apartats; aquest és el b), el del punt
    $P(-1,2,13)$.}
\end{errada}

\begin{errada}{Exercici 86 b)}
  \ondiu{$\vec{n}=(2,-2,1)$}
  \hadedir{$\vec{n}=(2,-2,-1)$}
  \motiu{El coeficient de $z$ del pla $2x-2y-z+1=0$ és $-1$. El mòdul, $3$,
    i la distància, $1$ u, no canvien.}
\end{errada}

\begin{errada}{Exercici 88, equació}
  \ondiu{$|\lambda+4|=\sqrt{6}$}
  \hadedir{$\dfrac{|\lambda+4|}{\sqrt{6}}=\sqrt{6}$}
  \motiu{Hi falta el denominador $\sqrt{6}$ de la fórmula de la distància.
    La línia següent, $|\lambda+4|=6$, ja és correcta.}
\end{errada}

\begin{errada}{Exercici 88, segon punt}
  \ondiu{$\lambda+4=-6 \;\rightarrow\; \lambda=-8 \;\rightarrow\; P(-6,19,-8)$}
  \hadedir{$\lambda+4=-6 \;\rightarrow\; \lambda=-10 \;\rightarrow\; P(-8,23,-10)$}
  \motiu{De $\lambda+4=-6$ surt $\lambda=-10$. El punt $(-6,19,-8)$ és a
    distància $\tfrac{4}{\sqrt{6}}$ del pla, no a $\sqrt{6}$.}
\end{errada}

%\senseerrades{76}
